\chapter{Games}
\newpage

\section{Suicide}
All throwers line up in a row. They all throw their boomerangs simultaneously on command. Those who catch their boomerangs are allowed to participate in the next round. Those who do not catch their own boomerangs are eliminated. The number of players is gradually reduced until one is left - the winner. Towards the end of the game, if the weather conditions are good, the game can be made more difficult for catching (e.g. by one-handed catching, catching behind the back, catching with the feet etc.) In order to ensure that there is little danger of injury, only small and light boomerangs should be used.

\section{Boomerang Rugby}
The rules for this game are very simple, making it even more fun. A player throws his boomerang and then he and as many other players as desired try to catch it. The one who catches, is thrower in the next round etc.

\section{Supercatch}
An MTA is thrown and while it is in the air, the thrower throws and catches a ``Fast-Catch'' boomerang as many times as possible. No catches are counted if the MTA boomerang is not caught. It is useful to have someone watch the MTA while the ``Fast-Catch'' boomerang is being thrown to advise the thrower when to start running for the MTA and where to catch it.
\newpage

\section{Fächerwurf}
The object of ``Fächerwurf'' is to throw three boomerangs in a row and then to catch them. The most difficult part is selecting the appropriate boomerangs and throwing them in the right order. There are many different strategies for success, two of which are:
Throw a boomerang with a short flight, then one with a medium flight and finally a boomerang with a longer flight.
If three similar boomerangs are thrown, they should hover slowly above the launching position and then gradually descend. They should be thrown at intervals of approximately 2 - 3 seconds.
The ``Fächerwurf'' can, of course, be played with more than 3 boomerangs.

\section{As many boomerangs as possible in the air}
In this game, try to keep as many boomerangs as possible simultaneously in the air. The boomerangs thrown before the first one touches the ground count towards the score.

\section{Schnelle Feile - Boomerang making race}
The goal is to make a boomerang as fast as possible out of Finnish birch using only a saw, a rasp and a file and then throw and catch it 5 times. The first one to finish wins.

\section{Throw off Fast Catch}
Two throwers launch simultaneously. The first one to catch his boomerang five times advances to the next round.
\newpage

\section{Team Relay}
The playing field consists of two concentric circles. The inner throwing circle is 4 metres and outer circle is 30 metres in  diameter. Two teams consisting of 3 - 6 throwers stand outside the 30 metre line. The 30 metre line is the start and finish line. On command, the first thrower of each team starts. The boomerang must be thrown from the 4 metre circle. After catching, the thrower touches the 4 metre circle before running back to the team. If he doesn't make the catch he returns to the 4 metre circle and throws again. Upon catching it, or recovering it if he does not catch it, he touches the 4 metre circle and runs back to tag a teammate at the 30 metre line to start his turn and so on. All throwers throw twice in rotation. The winning team is the one whose final thrower touches or crosses the 30 metre line first. 
To make the game a bit more interesting, try using only one boomerang for the team. This can, however, be a bit tricky for left-handers, as they will of course have difficulties throwing a right-handed boomerang.

\section{Chaos}
Two people (or more - in this example we will use two, but the rules are the same for more people) have one boomerang each, preferably the same style to insure fairness. Start the game by throwing your boomerang around the central marker. 

\subsection*{General Rules}
One point is scored each time a player catches his/her own boomerang-- a defensive point. Three Points are scored each time a player catches his/her opponent's boomerang--an offensive point. 

A throw may be made from any spot on the field, so long as it goes around the central marker. If it does not go around the marker, no defensive point may be scored, but offensive points may still be scored by an opponent catching it. There is no time limit, and players may throw their boomerangs again as soon as they catch or retrieve them. A winner is determined when one player reaches 11 points. 

\subsection*{Free Throws}
A Free Throw is when play has ceased and one player is allowed one throw without any interference from other players. If it is caught, the player gets a point, and normal play resumes. There are two circumstances for a Free Throw: (Starring Dick and Jane)

First, if there is a foul. Chaos is a ``non-contact'' sport like basketball, so if Dick slaps Jane's hand while trying to catch her boomerang, Jane may call "Foul!" She then gets a Free Throw. 

Second, if a player looses his/her boomerang (in a tree, over a fence, in tall grass, or just didn't watch where it landed), he can call ``Out of Play''. Play then stops and Jane gets ONE Free Throw while Dick looks for his lost boomerang. Calling ``Out of Play'' is in Dick's Best interest, even though Jane will probably get an easy point, because she will not be able to rack up point after point while Dick is combing through the neighbor's garden and climbing trees looking for his boomerang. Still, it is Dick's decision, whether it is worth it to call ``Out of Play'' or not. 

This game was originally invented by J and N Boomerangs, and rules was provided by Master-Designs.com